\documentclass{beamer}
\usepackage{graphicx}
\usepackage[utf8]{inputenc}
\usepackage[spanish]{babel}
\graphicspath{{Imagenes/}}
\usetheme{Dresden} 
%Otros temas: Antibes, AnnArbor, CambridgeUS, Berkeley, Goettingen, Bergen, Hannover, Dresden, Ilmenau, Berlin, Boadilla, Darmstadt

%Esto unicamente se utiliza para el tema Bergen
%\def\insertauthorindicator{¿Quién?}
%\def\insertdateindicator{Fecha}
%\title{Taller de herramientas computacionales}
%\author{Fernando Moctezuma Soto}
%\date{\today}

\title{Taller de herramientas computacionales}
\author{Fernando Moctezuma Soto}
\date{22/enero/2019}

\begin{document}
	
	\maketitle
	
	
	\begin{frame}
	\transblindshorizontal
		\frametitle{Primer presentación en LaTeX}
	\end{frame}


	\begin{center}
		\includegraphics[scale=0.50]{1.jpg}
	\end{center}

	\begin{frame}
		\frametitle{Segunda diapositiva}
		Esta es la segunda diapositiva.
	\end{frame}

	\begin{frame}[fragile]
		\frametitle{Tercera diapositiva}
		
		\begin{verbatim}
		#!/usr/bin/python2.7
		# -*- coding: utf-8 -*-
		
		'''
		Moctezuma Soto, Fernando
		310126070
		Taller de herramientas computacionales
		
		Programa que imprime "Hoy es miércoles"
		'''
		x = 10.5;  y = 1.0/3;  z = 15.3
		# x, y, z = 10.5, 1.0/3, 15.3
		
		H = """El punto en R3 es:
		(x, y, z) = (%.2f, %g, %G)
		""" % (x, y, z)
		print H
		
		G = """
		El punto en R3 es:
		(x, y, z) = ({laX:.2f}, {laY:g}, {laZ:G})
		""" .format(laX=x, laY=y, laZ=z)
		print G
		
		import math as m
		from math import sqrt 
		#from math import *
		
		x=16
		x=input("¿Cuál es el valor al que le quieres calcular la raiz?: ")
		print "La raíz cuadrada de %.2f es %f" % (x, m.sqrt(x))
		print sqrt(16.5)
		\end{verbatim}
		
	\end{frame}
		
\end{document}